\section{Ключевые принципы эффективного сканирования}
\label{sec:Chapter2} \index{Chapter2}

Прежде чем переходить к обзору конкретных сканеров и анализу их недостатков, имеет смысл рассмотреть, какие общие принципы лежат в основе работы современных сканеров уязвимостей и какие из этих принципов оказывают наибольшее влияние на эффективность сканирования.

\subsection{Модель сканирования образа}

Сканирование уязвимостей в OCI-образе сводится к решению следующих двух взаимосвязанных подзадач:

\begin{itemize}
\item \textit{индексирование (indexing)}~--- определение полного списка программных компонентов, содержащихся в образе, с указанием их имен и точных версий;
\item \textit{сопоставление с базой уязвимостей (matching)}~--- сравнение полученного списка с известными уязвимостями для соответствующих компонентов и формирование итогового отчёта.
\end{itemize}

Принципиальное различие между этими подзадачами заключается в том, что индексирование требует доступа к содержимому образа, в то время как сопоставление с базой уязвимостей оперирует уже готовым списком компонентов и от содержимого образа не зависит. Иными словами, индексирование~--- это анализ конкретного образа, тогда как сопоставление~--- это применение к результату индексирования общей для всех образов базы уязвимостей.

В качестве источников данных для сопоставления, как правило, используются базы уязвимостей, специфичные для типа компонента: для системных пакетов это дистрибутив-специфичные трекеры (Debian Security Tracker, Red Hat OVAL, Alpine SecDB и др.), для языковых пакетов~--- базы соответствующих экосистем (GitHub Advisory Database, RustSec, PyPI Advisory Database и др.), а в качестве общего агрегирующего источника~--- NVD~\cite{nvd}. Большинство сканеров объединяют эти источники в собственном формате, оптимизированном для быстрого поиска: у Trivy~\cite{trivy} это репозиторий \texttt{trivy-db}, у Grype~\cite{grype}~--- \texttt{grype-db}, у Clair~\cite{clair}~--- собственный формат, заполняемый отдельным компонентом. Согласованность этих источников между различными сканерами не гарантирована, поэтому результаты сканирования одного и того же образа разными сканерами могут заметно различаться, что является нормой и не считается ошибкой.

В большинстве современных сканеров (Trivy, Clair, Grype) описанные подзадачи концептуально или явно разделены. Так, Clair~\cite{clair} имеет два независимых HTTP-сервиса~--- indexer и matcher~--- каждый из которых занимается своей подзадачей. В системе Grype~\cite{grype} от компании Anchore разделение подведено к уровню инструментов: индексирование осуществляется отдельным инструментом Syft~\cite{syft}, тогда как сопоставление производит Grype. Trivy~\cite{trivy} объединяет обе подзадачи в рамках одной утилиты, однако внутри также реализует их раздельно.

\subsubsection{Преимущества разделения}

Разделение на индексирование и сопоставление имеет ряд практических преимуществ. Во-первых, результат индексирования можно сохранить и переиспользовать: появление новых записей в базе уязвимостей не требует повторного анализа образа, достаточно повторно применить процедуру сопоставления. Во-вторых, разделение позволяет выделить индексирование в отдельный этап CI-конвейера и сохранить её результат как самостоятельный артефакт; в этом случае хранение результата вместо самого образа существенно дешевле, что упрощает ретроспективный анализ. В-третьих, разделение в общем виде укладывается в концепцию \textit{SBOM (Software Bill of Materials)}, благодаря чему различные сканеры могут обмениваться результатами индексирования друг с другом.

\subsection{Слои как минимальная единица анализа}

При проведении индексирования необходимо определить, какую единицу содержимого считать минимальной для анализа. Хотя в конечном счёте сканер анализирует файловую систему контейнера, на практике принято использовать в качестве минимальной единицы анализа \textit{слой образа}. Такой выбор обусловлен несколькими соображениями.

Во-первых, как уже отмечалось в \hyperref[subsec:oci-structure]{разделе 2.2}, слои образа адресуются дайджестом своего содержимого. Это означает, что результат анализа слоя зависит исключительно от содержимого этого слоя и не зависит ни от других слоев, ни от метаданных образа в целом. Следовательно, результат анализа можно сохранить и повторно использовать при сканировании любого образа, в составе которого встречается тот же слой. Это особенно важно с учётом того, что слои активно повторно используются между различными образами: один и тот же базовый слой может встречаться в десятках или сотнях образов.

Во-вторых, использование слоя в качестве единицы анализа хорошо сочетается с инкрементальным обновлением образов в CI/CD-конвейерах: если при сборке нового образа изменился только верхний слой, то и заново анализировать необходимо только его, тогда как результаты анализа остальных слоев можно взять из кеша.

\subsubsection{Наложение слоев}

Анализ слоёв независимо друг от друга позволяет лишь получить набор частичных результатов~--- по одному на каждый слой. Чтобы по ним восстановить итоговый список компонентов образа, эти результаты необходимо совместить с учётом того, что файловая система контейнера получается \textit{последовательным наложением} слоёв в порядке, заданном манифестом.

Так, если один и тот же файл присутствует в нескольких слоях, итоговое содержимое определяется его версией в самом верхнем слое, поэтому при агрегации записи из более ранних слоёв должны быть вытеснены одноименными записями из более поздних. Аналогично whiteout-запись в верхнем слое \textit{логически удаляет} соответствующий файл из всех нижележащих слоёв, и его влияние должно быть исключено из общего результата индексирования~--- даже если в нижнем слое этот файл был фактически проанализирован.

Таким образом, агрегация результатов индексирования слоёв~--- это не просто их объединение, а воспроизведение тех же правил наложения, что применяются container runtime при сборке итоговой файловой системы контейнера. Корректное соблюдение этих правил обеспечивает соответствие результата индексирования фактическому содержимому образа.

\subsection{Роль метаданных в процессе сканирования}

Ключевая особенность процесса индексирования заключается в том, что информация о версиях программного обеспечения зачастую исходит не от самого этого ПО, а из \textit{метаданных} пакетных менеджеров и манифестов языковых зависимостей, присутствующих в образе. К таким метаданным относятся, например:

\begin{itemize}
\item \texttt{/var/lib/dpkg/*}~--- файл с информацией обо всех пакетах, установленных через apt/dpkg в Debian-based дистрибутивах;
\item \texttt{/lib/apk/db/installed}~--- аналогичный файл для пакетного менеджера apk в Alpine Linux;
\item \texttt{/var/lib/rpm/Packages}~--- база данных установленных пакетов в RPM-based дистрибутивах (CentOS, Fedora, Red Hat);
\item \texttt{requirements.txt}, \texttt{Pipfile.lock}, \texttt{poetry.lock} для Python-проектов;
\item \texttt{package.json}, \texttt{package-lock.json}, \texttt{yarn.lock} для Node.js-проектов;
\end{itemize}

Размеры перечисленных файлов метаданных, как правило, на много порядков меньше размеров самого образа: на типичный пакет приходится несколько сотен байт текстовой информации в соответствующем файле базы данных, а размер всего файла даже для образа с сотнями пакетов редко превышает пару мегабайт. Аналогично манифесты языковых зависимостей обычно занимают единицы или десятки килобайт. Иными словами, для сканера на уязвимости в подавляющем большинстве случаев \textit{интересна крайне малая часть содержимого образа}, в то время как остальные файлы~--- бинарные файлы программ, разделяемые библиотеки, файлы конфигурации, файлы данных~--- являются для сканера избыточными\footnote{Здесь выделяются исполняемые файлы Go-проектов~--- поскольку в образах они представлены в виде собранных исполняемых файлов, в котором уже скомпонованы все необходимые зависимости, то метаинформации, сохраняемой в отдельном манифесте, в их случае нет~--- вместо этого при компиляции бинарного файла Go-приложения метаинформация о версиях зависимостей \textit{вшивается в сам бинарный файл}. По этой причине при анализе зависимостей исполняемых файлов Go фактически необходимо запрашивать все исполняемые файлы в образе и анализировать представленные в них метаданные. Более подробно этот момент будет рассмотрен в следующих разделах.}.

Это наблюдение является ключевым для данной работы. Если бы сканер имел возможность загружать из registry только эту небольшую часть содержимого образа, скорость сканирования и нагрузка на инфраструктуру существенно бы сократились. Однако стандартный формат слоев OCI этого не предусматривает, поскольку слои сжаты алгоритмом gzip с непрерывной зависимостью между байтами: чтобы прочитать произвольный файл из слоя, необходимо полностью разжать поток до соответствующей позиции. \textbf{Этот недостаток и порождает основную проблему, исследуемую в данной работе}.

\newpage
